% !TeX root = ../../main.tex
\subsection{Model}

For guide $g$ in sequenced sample $i$, let $K_{gi}$ be its read count, let
$N_i$ be the full mapped-guide total for that sample, and let
$Y_{gi}=K_{gi}/N_i$ be the observed guide proportion.  Sample-level
quantities such as dose, time, treatment, donor, and batch are predictors even
when they are described biologically as experimental ``outcomes.''

The model is fitted separately for every guide, so the guide index is
suppressed below.  Conditional on a latent sample-specific abundance $P_i$,
\begin{align}
  K_i \mid P_i &\sim \operatorname{Binomial}(N_i,P_i),\\
  P_i &\sim \operatorname{Beta}(\mu_i\kappa,(1-\mu_i)\kappa).
\end{align}
The mean is linked to a row vector of sample covariates $\mathbf{x}_i^\top$ by
\begin{equation}
  \logit(\mu_i)
  = \log\left(\frac{\mu_i}{1-\mu_i}\right)
  = \mathbf{x}_i^\top\boldsymbol{\beta}.
  \label{eq:mean}
\end{equation}
The design can contain an intercept, a continuous phenotype, batch indicators,
interactions, spline terms, and other prespecified adjustment variables.

It is useful to express the beta precision $\kappa$ as an intraclass
correlation
\begin{equation}
  \rho = \frac{1}{\kappa+1}, \qquad 0\leq\rho<1.
\end{equation}
Marginally,
\begin{align}
  \E(K_i) &= N_i\mu_i,\\
  \Var(K_i) &=
  N_i\mu_i(1-\mu_i)\{1+(N_i-1)\rho\}, \label{eq:countvar}\\
  \Var(Y_i) &=
  \mu_i(1-\mu_i)
  \left\{\frac{1-\rho}{N_i}+\rho\right\}. \label{eq:propvar}
\end{align}
Equation~\eqref{eq:propvar} cleanly separates sequencing uncertainty, which
decreases with $N_i$, from between-sample heterogeneity, which does not vanish
when sequencing depth becomes arbitrarily large.

\subsubsection{Interpretation of the continuous coefficient}

Suppose the model contains a continuous phenotype $d_i$ with coefficient
$\beta_d$.  Then $\exp(\beta_d)$ is the multiplicative change in the odds of
observing that guide for a one-unit increase in $d_i$, conditional on the other
covariates.  Standardizing $d_i$ before fitting makes $\exp(\beta_d)$ the odds
ratio per sample standard deviation and makes guide effects directly
comparable.  For rare guides, odds and proportions are similar, so $\beta_d$ is
also approximately a log abundance ratio per unit of $d_i$.
